%% Magic2 — Nature Methods Brief Communication
%% Draft — April 2026
%% Benchmark section is a placeholder pending final numbers.

\documentclass[11pt]{article}
\pdfoutput=1

\usepackage[margin=2.5cm]{geometry}
\usepackage{times}
\usepackage{graphicx}
\usepackage{booktabs}
\usepackage{hyperref}
\usepackage{natbib}
\usepackage{authblk}
\usepackage{amsmath}
\usepackage{microtype}

\hypersetup{colorlinks=true, linkcolor=blue, citecolor=blue, urlcolor=blue}

%% ---------------------------------------------------------------
\title{\textbf{Magic2: a unified, self-configuring aligner for short-
       and long-read RNA-seq with comprehensive quality-control output}}

\author[1]{Jean Thierry-Mieg}
\author[1]{Danielle Thierry-Mieg}
\author[1]{Greg Boratyn}
\affil[1]{National Center for Biotechnology Information, National Library
          of Medicine, National Institutes of Health, Bethesda, MD, USA}

\date{}
%% ---------------------------------------------------------------

\begin{document}
\maketitle

\begin{abstract}
We present Magic2, a fast, self-configuring RNA-seq and DNA-seq aligner
that operates identically on short (Illumina) and long (Nanopore,
PacBio) reads without requiring the user to specify read type, adapter
sequences, strandedness, or the number of threads.
In a single pass, Magic2 produces alignments, splice-junction tables,
gene-level expression counts, strand-specific coverage tracks,
poly(A)/poly(T) tail calls, and a rich set of quality-control metrics
— all directly comparable across runs and laboratories.
Benchmarks on 35 public datasets (totalling $\sim$70\,Gb, including
human, chimpanzee, and lemur reads aligned to the human T2T reference,
as well as Nanopore and PacBio long-read datasets) show that Magic2
matches or exceeds leading aligners (STAR, HISAT2, Minimap2,
Magic-BLAST) on alignment rate, concordant-pair fraction, mismatch
rate, and transcript-level accuracy, while achieving competitive
elapsed time and substantially lower CPU time than STAR.
Magic2 is freely available at
\url{https://github.com/NLM-DIR/Magic2}.
\end{abstract}

%% ---------------------------------------------------------------
\section*{Introduction}

High-throughput RNA sequencing is now the dominant approach for
measuring gene expression, yet the computational chain that transforms
raw reads into biological conclusions remains fragile and poorly
standardised.
Small differences in aligner choice, parameter settings, or
post-processing steps routinely produce divergent results from
identical input data~\citep{maqc2014,seqc2019}.
This fragility has driven the impractical recommendation that authors
record every parameter down to the exact version of each software
component — an implicit acknowledgement that reproducibility is not
yet guaranteed by design.

A contributing source of variability is the proliferation of
specialised tools: STAR~\citep{dobin2013} and
HISAT2~\citep{kim2019} are optimised for short Illumina reads,
while Minimap2~\citep{li2018} targets long, error-prone reads.
Users must master multiple programs with different parameter
conventions and output formats, and the choice of default parameters
can unintentionally bias benchmark comparisons.
No single aligner currently handles the full spectrum of read
technologies with consistent accuracy and without manual
reconfiguration.

Magic2 addresses these limitations through three design principles.
First, full automation: read length, paired-end geometry, adapter
sequences, and sequencing platform are inferred from the data;
internal parameters such as seed length and maximum intron size are
derived from genome size.
Second, single-pass richness: alignments, expression counts,
splice-junction support, coverage tracks, poly(A) calls, and QC
metrics are all produced in one run, eliminating downstream
processing tools.
Third, standardised QC: every run emits an identical set of quality
metrics (fraction of reads mapping to CDS, introns, intergenic
regions; strand-specificity score; concordant-pair rate; mismatch
distribution), enabling large-scale, cross-laboratory comparisons of
SRA datasets.

%% ---------------------------------------------------------------
\section*{Results}

\subsection*{Overview of Magic2}

Magic2 accepts a reference genome (FASTA) and an optional annotation
(GTF/GFF) to build an on-disk seed index.
At run time, one or more SRA accession numbers or local FASTA/FASTQ
files are supplied; no other configuration is required.
Magic2 detects read length, adapter sequences, paired-end status, and
strandedness automatically via a dry pass over the first reads.
The same executable is used for Illumina short reads, PacBio CCS, and
Nanopore long reads.

Internally, Magic2 follows a sort–merge paradigm rather than hash-table
lookup: seed records are generated for reads and genome in large
sequential blocks, sorted with a cache-efficient radix sort (optionally
GPU-accelerated via CUDA Thrust), and joined by simultaneous linear
traversal of both sorted lists.
This design keeps the dominant computation in fast CPU cache rather than
in random main-memory accesses, and scales naturally across NUMA
architectures without user-supplied thread counts
(see Supplementary Material for full details).

The pipeline is structured as concurrent agents communicating through
bounded, lock-free channels.
Agents process independent data blocks in parallel; backpressure and
channel closure guarantee deadlock-free execution and no data loss.

\subsection*{Outputs}

In a single pass Magic2 produces:
(i) primary alignments in BAM format (or the compact \texttt{.hits}
    format with \texttt{--hits});
(ii) splice-junction tables with per-junction read support;
(iii) gene-level expression counts with bidirectional-overlap resolution;
(iv) strand-specific wiggle coverage tracks (plus read-start/end profiles
     for transcript boundary detection);
(v) poly(A)/poly(T) tail calls and trans-spliced leader identification
    (in \textit{C.\ elegans});
(vi) adapter consensus sequences (de-novo or from a supplied list);
(vii) a standardised QC table in Tag–Sample Format (TSF) covering
     alignment rates, mapping-class fractions, strand-specificity,
     mismatch distributions, and concordant-pair statistics.

\subsection*{Benchmark}

\textbf{[PLACEHOLDER — to be completed with final numbers.]}

\medskip
\noindent
We benchmarked Magic2 against STAR (v.\,X.X), HISAT2 (v.\,X.X),
Minimap2 (v.\,X.X), and Magic-BLAST (v.\,X.X) on a panel of 35 publicly
available datasets downloaded from the NCBI SRA, each comprising
approximately 2\,Gb of reads (Table~\ref{tab:datasets}).
The panel includes:
\begin{itemize}
  \item single-end and paired-end Illumina 150\,nt reads (majority human);
  \item six cross-species datasets (chimpanzee to lemur) aligned to the
        human T2T reference genome, providing a controlled divergence
        gradient for sensitivity testing;
  \item approximately 14 Nanopore and PacBio long-read datasets;
  \item RefSeq mRNA sequences used as a zero-SNP, perfect-alignment
        control.
\end{itemize}
All tools were run with default parameters except where noted; adapter
trimming for competing tools used Trimmomatic (v.\,X.X) with
manufacturer-recommended settings.
Magic2 required no additional configuration for any dataset.

\medskip
\noindent
\textbf{[INSERT summary figure: heatmap or radar plot of per-tool,
per-metric scores across all 35 datasets.]}

\medskip
\noindent
\textbf{[INSERT Table 2: CPU time, elapsed time, peak memory for each
tool on the human paired-end reference dataset.]}

\medskip
\noindent
Across the full panel, Magic2 ranked first or tied for first on
\textbf{[X/Y]} quality metrics (alignment rate, concordant-pair
fraction, perfect-alignment count, mismatch rate, and
transcript-level precision).
HISAT2 was the fastest tool in elapsed time but showed substantially
lower sensitivity, particularly on divergent-species and long-read
datasets.
On the RefSeq control, Magic2 achieved \textbf{[X]}\,\% perfect
alignments, confirming correctness of the seed–sort–merge pipeline.
CPU time for Magic2 was \textbf{[X]}-fold lower than STAR and
comparable to Minimap2 on 16 threads.

%% ---------------------------------------------------------------
\section*{Discussion}

Magic2 achieves the primary goal set out in its design: a single
executable that handles short and long reads, requires no
configuration, and produces a richer set of outputs than any
comparable tool in a single pass.
The benchmark results confirm that automation of parameter selection
does not come at the cost of accuracy — Magic2 matches or exceeds
specialised tools on their own preferred data types.

The standardised QC output is, we argue, the feature with the broadest
long-term impact.
Applied at scale to the SRA, it would allow independent assessment of
data quality without rerunning every published analysis, directly
addressing the reproducibility concerns raised in the Introduction.

Magic2 is under active development.
Planned near-term additions include GitHub Actions CI/CD for automated
multi-platform builds (Linux x86-64, Linux ARM, macOS), a regression
test suite, and expanded documentation.

%% ---------------------------------------------------------------
\section*{Methods}

\subsection*{Index construction}

Given a reference genome in FASTA format and an optional GTF/GFF
annotation, Magic2 builds an on-disk seed index.
The seed length $k$ (default 16 for genomes $<$1\,Gb, 18 for
human/mouse) is selected automatically from genome size.
Seeds are sub-sampled with a rolling-min-hash approach (window $p$,
default 5) to achieve an average inter-seed spacing of $p/2$ bases
with a guaranteed maximum gap of $p$.
Seeds with genomic multiplicity $> Q$ (default 81) are discarded to
limit hits from repetitive elements.
When annotation is provided, junction-spanning seeds from all annotated
exon--exon junctions are added to the index to improve detection of
reads crossing short exons.
Full details are given in Supplementary Section~2.

\subsection*{Alignment}

Read seeds are generated using the same parameters as the genome index,
merged-sorted with the genomic seed table, and joined by simultaneous
linear traversal.
Seed matches are extended to full alignments using a banded
Smith–Waterman variant; intron boundaries are refined by splice-signal
scoring (GT-AG, AT-AC, and non-canonical).
Poly(A)/poly(T) tails and trans-spliced leader sequences are detected
during the same pass.
See Supplementary Sections~3--5 for algorithmic details.

\subsection*{Parallelism}

Magic2 spawns concurrent agents connected by bounded lock-free channels.
On machines with multiple NUMA nodes, it re-spawns itself via
\texttt{numactl} to bind to the least-loaded node; this step is skipped
automatically in containers and on single-node machines.
Users do not set a thread count; Magic2 manages concurrency
automatically.

\subsection*{Benchmark datasets and evaluation}

[PLACEHOLDER — dataset accession table and evaluation metric
definitions to be inserted here.]

%% ---------------------------------------------------------------
\section*{Code availability}

Magic2 source code is available at
\url{https://github.com/NLM-DIR/Magic2} (branch \texttt{master})
under [LICENSE].
Pre-built binaries for Linux x86-64 are attached to each GitHub
Release.

%% ---------------------------------------------------------------
\section*{Acknowledgements}

[To be completed.]

%% ---------------------------------------------------------------
\bibliographystyle{naturemag}
\bibliography{magic2_refs}

%% ---------------------------------------------------------------
\clearpage
\appendix

\section*{Supplementary Material}
See separate file \texttt{magic2\_supplementary.tex}.

\end{document}
